% ----- CHAUPAI 20 -----

\newpage

\subsection*{\deva{विंशतितम चौपाई} \(Twentieth Chaupai\)}
\addcontentsline{toc}{subsection}{\deva{विंशतितम चौपाई} \(Twentieth Chaupai\) — \deva{दुर्गम काज...}}

\begin{minipage}[t]{0.55\textwidth}
\vspace{0pt}

{\large \linespread{1.5}\selectfont
\deva{दुर्गम काज जगत के जेते।}\\
\deva{सुगम अनुग्रह तुम्हरे तेते॥}
\par}

\vspace{0.5em}

{\large \linespread{1.5}\selectfont
\telu{దుర్గమ కాజ జగత కే జేతే।}\\
\telu{సుగమ అనుగ్రహ తుమ్హరే తేతే॥}
\par}

\vspace{0.5em}

{\large \linespread{1.3}\selectfont
\textit{durgama kāja jagata ke jete |}\\
\textit{sugama anugraha tumhare tete ||}
\par}
\end{minipage}%
\hfill
\begin{minipage}[t]{0.42\textwidth}
\vspace{0pt}
\begin{tcolorbox}[anchorbox]
\textbf{The Ultimate Problem Solver:} He takes 'impossible' tasks (\textit{Durgam}) and makes them manageable (\textit{Sugam}). When students face overwhelming exams, massive projects, or complex challenges, remembering his laser-like focus and unmatched work ethic provides a blueprint for making the impossible solvable.
\vspace{0.5em}\newline He specialized in taking wildly impossible tasks (Durgama) and breaking them down into solvable, manageable steps.\newline\null\hfill\href{https://www.valmiki.iitk.ac.in/sloka?field_kanda_tid=5&language=dv&field_sarga_value=39}{\textit{Sundara Kanda, 39:53}}
\end{tcolorbox}
\end{minipage}

\vspace{0.5em}
\subsubsection*{\textbf{Visual Rhythm Map:}}
\begin{table}[h!]
\centering
\renewcommand{\arraystretch}{1.2}
\setlength{\tabcolsep}{0pt}
\begin{tabularx}{\textwidth}{|*{16}{>{\centering\arraybackslash\hsize=1.0\hsize}X|}}
\hline
\multicolumn{4}{|>{\centering\arraybackslash\hsize=4.0\hsize}X|}{\textbf{\guru{दु}\laghu{र्ग}\laghu{म}} \par (4)} &
\multicolumn{3}{>{\centering\arraybackslash\hsize=3.0\hsize}X|}{\textbf{\guru{का}\laghu{ज}} \par (3)} &
\multicolumn{3}{>{\centering\arraybackslash\hsize=3.0\hsize}X|}{\textbf{\laghu{ज}\laghu{ग}\laghu{त}} \par (3)} &
\multicolumn{2}{>{\centering\arraybackslash\hsize=2.0\hsize}X|}{\textbf{\guru{के}} \par (2)} &
\multicolumn{4}{>{\centering\arraybackslash\hsize=4.0\hsize}X|}{\textbf{\guru{जे}\guru{ते}} \par (4)} \\
\hline
\end{tabularx}
\vspace{-1.5em}
\end{table}

\begin{table}[h!]
\centering
\renewcommand{\arraystretch}{1.2}
\setlength{\tabcolsep}{0pt}
\begin{tabularx}{\textwidth}{|*{16}{>{\centering\arraybackslash\hsize=1.0\hsize}X|}}
\hline
\multicolumn{3}{|>{\centering\arraybackslash\hsize=3.0\hsize}X|}{\textbf{\laghu{सु}\laghu{ग}\laghu{म}} \par (3)} &
\multicolumn{5}{>{\centering\arraybackslash\hsize=5.0\hsize}X|}{\textbf{\laghu{अ}\guru{नु}\laghu{ग्र}\laghu{ह}} \par (5)} &
\multicolumn{4}{>{\centering\arraybackslash\hsize=4.0\hsize}X|}{\textbf{\laghu{तु}\laghu{म्ह}\guru{रे}} \par (4)} &
\multicolumn{4}{>{\centering\arraybackslash\hsize=4.0\hsize}X|}{\textbf{\guru{ते}\guru{ते}} \par (4)} \\
\hline
\end{tabularx}
\vspace{-1.5em}
\end{table}

\vspace{0.5em}
\textbf{ \deva{सरल-संस्कृतम्}:}\\
 \deva{जगतः यावन्ति अपि दुर्गमाणि (कठिनानि) कार्याणि सन्ति}\\
 \deva{तानि सर्वाणि भवतः अनुग्रहेण (कृपया) सुगमानि भवन्ति॥}\\

Whatever difficult tasks exist in the world\\
they all become easily attainable by your grace.


\vspace{1em}
\newpage
\textbf{\deva{प्रतिपदार्थः} (Meaning of each word):}
\begin{table}[h!]
\centering
\renewcommand{\arraystretch}{1.2}
\begin{tabularx}{\textwidth}{@{} l l X X @{}}
\toprule
\textbf{अवधी} & \textbf{संस्कृतम्} & \textbf{English} & \textbf{Grammar \& Notes} \\
\midrule
\deva{दुर्गम} / \telu{దుర్గమ} / \textit{durgama}  &  \deva{दुर्गमानि}  &  Difficult / Hard to cross  &   \\
\deva{काज} / \telu{కాజ} / \textit{kāja}  &  \deva{कार्याणि}  &  Tasks / Work  & \\ 
\deva{जगत के} / \telu{జగత కే} / \textit{jagata ke}  &  \deva{जगतः}  &  Of the world  & Genitive Postposition `ke` \\
\deva{जेते} / \telu{జేతే} / \textit{jete}  &  \deva{यावन्ति}  &  However many  & Rhymes with `tete` \\
\deva{सुगम} / \telu{సుగమ} / \textit{sugama}  &  \deva{सुగమాని}  &  Accessible / Easy  &   \\
\deva{अनुग्रह} / \telu{అనుగ్రహ} / \textit{anugraha}  &  \deva{अनुग्रहेण}  &  By grace / blessing  &   \\
\deva{तुम्हरे} \gotchaicon / \telu{తుమ్హరే} / \textit{tumhare}  &  \deva{భవతః}  &  Your  &   \\
\deva{तेते} / \telu{తేతే} / \textit{tete}  &  \deva{तावन्ति}  &  That many  &   \\
\bottomrule
\end{tabularx}
\end{table}

\begin{tcolorbox}[gotchabox]
\begin{itemize}
    \item \textbf{Aspirated Consonants (\deva{म्ह} and \deva{न्ह}):} In words like \deva{तुम्हरे} (Tumhare) and \deva{कीन्हा} (Keenhaa), the `h` merges with `m` and `n` to form a single aspirated sound. In classical Awadhi meter, this does \textit{not} make the preceding short vowel heavy. Therefore, \deva{तु} (tu) in Tumhare remains exactly 1 beat!
\end{itemize}
\end{tcolorbox}

\newpage
